\section{Proof of main theorem}
Here we present the proof of Theorem \ref{thm:looijenga}. Recall that $X$ is a path connected topological space with basepoints $a,b$. Our strategy will be to consider the spectral sequence associated to the two-sided cobar construction. We will show that it converges to the group $H_n(X^n, X(n)^a_b)$ and conclude by the previous discussion relating the cobar construction to the fundamental group. 

\subsection{Building the spectral sequence} Recall the cosimplicial path space $P_{a,b}^\bullet X$ of Example \ref{example:cosimplicial_path_space}. From now we denote it by $P^\bullet$ for simplicity. Notice first that the union of the images of the coface maps $d^i: X^{n-1}\to X^n$ is exactly the subspace $X(n)^a_b$. This is a good sign.

Now consider the bicomplex $\mathscr{C}_{\bullet,\bullet}$ which is obtained by applying singular chains to $P^\bullet$: 
\[\mathscr{C}_{p,q} = C_{-p}(P^q)\]
with ``horizontal'' differential $d^\mathscr{C}_H: C_p(P^q)\to C_{p+1}(P^q)$ given by the singular boundary map, and ``vertical'' differential $d^V_\mathscr{C}: C_p(P^q) \to C_p(P^{q+1})$ given by the differential of the Moore complex of the cosimplicial abelian group $C_p(P^\bullet)$. 

% \begin{prop}
%  $C^\bullet_\bullet$ is a bicomplex.
%  \begin{proof}
%  That $d^\mathscr{C}_H$ and $d^V_\mathscr{C}$ both square to zero follows from the fact that each row is a chain complex and each row is a cochain complex. It remains to check commutativity of the squares: 
 
% \[\begin{tikzcd}[cramped]
%     % https://q.uiver.app/#q=WzAsNCxbMSwxLCJDX3AoWF5xKSJdLFsxLDAsIkNfcChYXntxKzF9KSJdLFswLDAsIkNfe3AtMX0oWF57cSsxfSkiXSxbMCwxLCJDX3twLTF9KFhee3F9KSJdLFswLDFdLFswLDNdLFszLDJdLFsxLDJdXQ==
% 	{C_{p-1}(X^{q+1})} & {C_p(X^{q+1})} \\
% 	{C_{p-1}(X^{q})} & {C_p(X^q)}
% 	\arrow[from=1-2, to=1-1]
% 	\arrow[from=2-1, to=1-1]
% 	\arrow[from=2-2, to=1-2]
% 	\arrow[from=2-2, to=2-1]
% \end{tikzcd}\]
% Starting with a map $f: \Delta^p \to X^q$, if we first go up then left, we get the map 
% \[ \Delta^{p-1}\xrightarrow{([p]\to [p-1])^*}\Delta^p \xrightarrow{f} X^q \xrightarrow{\sum (-1)^id^i} X^{q+1}\]
% and proceeding in the other direction gives the exact same map.
% (This is really obvious, might not need to include).
%  \end{proof}
% \end{prop}

Each column $p$ of this bicomplex forms a cochain complex $\mathscr{C}_{p,\bullet}$, the Moore complex of $C_p(F^\bullet)$. Consider now the bicomplex $\mathscr{N}_{\bullet,\bullet}$ whose columns are instead the normalized cochain complexes of $C_p(F^\bullet)$. The map $\mathscr{C}\twoheadrightarrow \mathscr{N}$ is a quasi-isomorphism of bicomplexes by Dold--Kan and Lemma \ref{lem:quasiiso_bicomplex}. Let us spell out what the groups in $\mathscr{N}$ actually are. Following the definition, we have 
\[\mathscr{N}_{p,q} = \coker\bigoplus_{i=1}^{n} ((d^i)_*: C_p(X^{n-1})\to C^p(X^n)) = C^p(X^n) / \lb \sum_{i=1}^n \im (d^i)_*\rb  = C_p(X^n, X(n)_b).\]
Moreover, the vertical differential $\mathscr{N}_{p,q} \to \mathscr{N}_{p,q+1}$ is induced by the zeroth coface map $d^0$. Consider now the bicomplex obtained from $\mathscr{N}$ by setting all rows above $q=n$ to zero, which we will denote by $(\mathscr{N}^{\leq n})^\bullet_\bullet$.
\begin{prop}
    Filtering $(\mathscr{N}^{\leq n})$ by the lower index, we obtain a second-quadrant spectral sequence collapsing on the second page with 
    \[E^2_{-n,n} = H_n(X^n, X(n)^a_b), \quad E^2_{0,0} = \begin{cases}
        \Z & a=b\\
        0 &  a\neq b
    \end{cases}\]
    converging to $H_\bullet\tot(\mathscr{N}^{\leq n})$.
    \begin{proof}
        The filtration is bounded below and exhaustive, so the spectral sequence converges. To compute the $E^2$-page, let us begin on $E^0$, whose $-p$th column is the following normalized cochain complex (beginning below in degree $q=0$):
        \[C_p(*)\to C_p(X, \{b\})\to \cdots\to C_p(X^n, X(n)_b)\to 0 \to \cdots\]
        of the Moore complex 
        \begin{equation}\label{eqn:moore_complex}
            C_p(*)\to C_p(X) \to \cdots \to C_p(X^n)\to 0 \to \cdots
        \end{equation}
        so they have the same cohomology. But now (\ref{eqn:moore_complex}) is just the cochain complex obtained by taking the free abelian group on the cosimplicial set $[q] \mapsto \Map(\Delta^p, F^q)$, with terms $q>n$ set to zero. Therefore, by Proposition \ref{prop:goodwillie_cor}, we have that $E^1_{-p,q} = 0$ for $q \neq 0,n$. For $q=n$ and $q=0$, we have 
        \[E^1_{-p,n} = C_p(X^n, X(n)_b) /\im(d^0) = C_p(X^n, X(n)^a_b)\]
        \[E^1_{-p,0} = \ker(d^0: C_p(*)\to C_p(X, \{b\})) = \begin{cases}
        \Z & p=0 , a=b \\
        0 & \text{otherwise}
        \end{cases}.\]

        On the $E^1$-page, the differential is now induced by singular boundary map. On the first nonzero row $q=n$ of the $E^1$-page, we just take singular homology to obtain
        \[E^2_{-p, n} = H_p(X^n, X(n)^a_b).\]
        And on the other possibly nonzero row $q=0$, we have 
        \[E^2_{-p, 0} = \begin{cases}
            H_p(\Z\leftarrow 0 \leftarrow \cdots) & a=b \\
            0 & a\neq b
        \end{cases}\]
        as desired. 

        At this point, because the only nonzero rows are $q=n$ and possibly $q=0$, the differentials on the pages $E^{\geq 2}$ are all zero, except possibly on page $E^{n+1}$, where we have a differential $d_{n+1}$ of bidegree $(n+1, -n)$ from $E^{n+1}_{-n-1, n} = H_{n+1}(X^n, X(n)^a_b)$ to $E^{n+1}_{0,0} = \Z$. To see that $d_{n+1} = 0$, consider the map of pairs $f: (X,x)\to (*,*)$. Repeating the constructions in this section, we get a spectral sequence $D_{\bullet,\bullet}$ computing $H_n(*^n, *(n)^*_*) = 0$. By functoriality, $f$ induces a morphism of spectral sequences. Specializing to the $n+1$th page, we have a commutative square 
		
		\[\begin{tikzcd}[cramped]
			% https://q.uiver.app/#q=WzAsOCxbMSwwLCJFXntuKzF9X3tuKzEsbn0iXSxbMSwxLCJEXntuKzF9X3tuKzEsbn0iXSxbMiwwLCJFXntuKzF9X3swLDB9Il0sWzIsMSwiRF57bisxfV97MCwwfSJdLFswLDAsIkhfe24rMX0oWF5uLFgobileeF94KSJdLFszLDAsIlxcWiJdLFswLDEsIkhfe24rMX0oKiwqKT0wIl0sWzMsMSwiXFxaIl0sWzAsMSwiZl8qIiwyXSxbMCwyLCJkXkVfe24rMX0iXSxbMSwzLCJkXkRfe24rMX0iLDJdLFsyLDMsImZfKiJdLFs0LDAsIiIsMix7InN0eWxlIjp7ImhlYWQiOnsibmFtZSI6Im5vbmUifX19XSxbNiwxLCIiLDIseyJzdHlsZSI6eyJoZWFkIjp7Im5hbWUiOiJub25lIn19fV0sWzIsNSwiIiwyLHsic3R5bGUiOnsiaGVhZCI6eyJuYW1lIjoibm9uZSJ9fX1dLFszLDcsIiIsMix7InN0eWxlIjp7ImhlYWQiOnsibmFtZSI6Im5vbmUifX19XV0=
			{H_{n+1}(X^n,X(n)^x_x)} & {E^{n+1}_{-n-1,n}} & {E^{n+1}_{0,0}} & \Z \\
			{H_{n+1}(*,*)=0} & {D^{n+1}_{-n-1,n}} & {D^{n+1}_{0,0}} & \Z
			\arrow[equal, from=1-1, to=1-2]
			\arrow["{d^E_{n+1}}", from=1-2, to=1-3]
			\arrow["{f_*}"', from=1-2, to=2-2]
			\arrow[equal, from=1-3, to=1-4]
			\arrow["{f_*}", from=1-3, to=2-3]
			\arrow[equal, from=2-1, to=2-2]
			\arrow["{d^D_{n+1}}"', from=2-2, to=2-3]
			\arrow[equal, from=2-3, to=2-4]
		\end{tikzcd}\]
		But now the $f_*: \Z\to\Z$ is the identity since it is the induced map $H_0(X,x)\to H_0(*,*)$, forcing $d^E_{n+1} = 0$, so $E^2 = E^{\infty}$ as desired.
    \end{proof}
\end{prop}
\begin{cor}\label{cor:spectral_convergence}
We have a surjection 
\[H_0(\tot(\mathscr{N}^{\leq n})) \twoheadrightarrow H_n(X^n, X(n)^a_b)\]
whose kernel is $\Z$ if $a=b$, and trivial otherwise.
\begin{proof}
    Follows from convergence of the spectral sequence.
\end{proof}
\end{cor}
\subsection{Functoriality}
It remains to compute $H_0$ of this totalization. For this we will return to our original bicomplex $\mathscr{C}$ and show that its truncated totalization is quasi-isomorphic to the totalization of the two-sided cobar bicomplex. Instead of normalizing the columns, let us first consider the rows $\mathscr{C}_{\bullet,q}$, which are the singular chain complexes $C_\bullet(X^q)$. Then we can apply the topological Alexander--Whitney map on each row:
\[\mathscr{C}_{-\bullet,q} = C_\bullet(X^q) \to C_\bullet(X)^{\tensor q}\]
to obtain a new bicomplex, which we will call $\mathscr{D}$. By definition of tensor product of chain complexes, $\mathscr{D}$ has a bigrading given by
\begin{equation}
\mathscr{D}_{-p,q} = (C_\bullet(X)^{\tensor q})_p = \bigoplus_{n_1 + \cdots + n_q = p} C_{n_1}(X) \tensor \cdots \tensor C_{n_q}(X).
\end{equation}
Let $\mathscr{C}^{\leq n}$ and $\mathscr{D}^{\leq n}$ be the bicomplexes obtained from $\mathscr{C}$ and $\mathscr{D}$ by setting all rows of above $q=n$ to zero. 

\begin{lem}\label{lem:quasi_iso_CtoD}
	The map $\mathscr{C}^{\leq n} \to \mathscr{D}^{\leq n}$ induced by Alexander--Whitney is a quasi-isomorphism of bicomplexes, with inverse quasi-isomorphism induced by the Eilenberg--Zilber map.
	\begin{proof}
		Combine Corollary \ref{cor:topological_aw}, Theorem \ref{thm:eilenberg_zilber}, and Lemma \ref{lem:quasiiso_bicomplex},
	\end{proof}
\end{lem}
\begin{thm}\label{thm:two_sided_cobar_D}
	Let $C_\bullet(-)$ be the singular chains functor. Regard $C_\bullet(\{a\})$ as a right $C_\bullet(X)$-comodule via the right comodule action 
	\[\Delta^R: C_\bullet(\{a\})\xrightarrow{\AW} C_\bullet(\{a\})\tensor C_\bullet(\{a\}) \xrightarrow{\id \tensor \iota_a} C_\bullet(\{a\}) \tensor C_\bullet(X)\]
	and likewise for $C_\bullet(\{b\})$ as a left $C_\bullet(X)$-comodule. Then there is a quasi-isomorphism 
	\[\Cobar^{\leq n}(C_\bullet(\{b\}), C_\bullet(X), C_\bullet(\{a\})) \to \mathscr{D}^{\leq n}.\]
	\begin{proof}
		Denote $\Cobar(\cdots)$ by $\Cobar$. By flatness, Lemma \ref{lem:graded_quism}, and Lemma \ref{lem:normalized_inclusion_is_graded_quism}, it suffices to exhibit a quasi-isomorphism from 
		\begin{equation}\label{eqn:two_sided_normalized_cobar}
		N\Cobar^{\leq n} : = \Cobar^{\leq n}(NC_\bullet(\{b\}), NC_\bullet(X), NC_\bullet(\{a\}))
		\end{equation}
		to $\mathscr{D}^{\leq n}$. Now by Lemma \ref{lem:quasiiso_bicomplex}, it suffices to exhibit a quasi-isomorphism from (\ref{eqn:two_sided_normalized_cobar}) to $N\mathscr{D}^{\leq n}$, the truncation of the bicomplex obtained by replacing each column of $\mathscr{C}$ with its normalized chains, then applying the Alexander--Whitney map, as in the diagram below:
		\[\begin{tikzcd}[cramped]
			% https://q.uiver.app/#q=WzAsNCxbMCwwLCJDX1xcYnVsbGV0KFhecSkiXSxbMCwxLCJOQ19cXGJ1bGxldChYXnEpIl0sWzEsMCwiQ19cXGJ1bGxldChYKV57XFx0ZW5zb3IgcX0gPSBcXG1hdGhzY3J7RH1fXFxidWxsZXRecSJdLFsxLDEsIk5DX1xcYnVsbGV0KFgpXntcXHRlbnNvciBxfSA9Ok5cXG1hdGhzY3J7RH1ecV9cXGJ1bGxldCJdLFswLDIsIlxcQVdfcSJdLFsxLDMsIlxcQVdfcSIsMl0sWzEsMCwiIiwwLHsic3R5bGUiOnsidGFpbCI6eyJuYW1lIjoiaG9vayIsInNpZGUiOiJ0b3AifX19XSxbMywyLCIiLDIseyJzdHlsZSI6eyJ0YWlsIjp7Im5hbWUiOiJob29rIiwic2lkZSI6InRvcCJ9fX1dXQ==
			{C_\bullet(X^q)} & {C_\bullet(X)^{\tensor q} = \mathscr{D}_{\bullet,q}} \\
			{NC_\bullet(X^q)} & {NC_\bullet(X)^{\tensor q} =:N\mathscr{D}_{\bullet, q}}
			\arrow["{\AW_q}", from=1-1, to=1-2]
			\arrow[hook, from=2-1, to=1-1]
			\arrow["{\AW_q}"', from=2-1, to=2-2]
			\arrow[hook, from=2-2, to=1-2]
		\end{tikzcd}\]
		Notice that $NC_\bullet(\{a\})$ and $NC_\bullet(\{b\})$ are just the chain complexes with a $\Z$ term concentrated in degree zero, i.e. the identity object in the symmetric monoidal structure on $\catname{Ch}_+(\catname{Ab})$. Hence we have natural isomorphisms 
		\[\psi_a: NC_\bullet(\{a\}) \tensor NC_\bullet(X) \to NC_\bullet(X)\]
		\[\psi_b: NC_\bullet(X) \tensor NC_\bullet(\{b\}) \to NC_\bullet(X)\]
		which combine to give a natural isomorphism $\phi:= \psi_a \tensor \id^{\tensor \bullet}\tensor \psi_b$:
		\begin{equation}
			\begin{split}
			\phi: N\Cobar_{-p,q} &=  \bigoplus_{n + b_1 + \cdots + b_q + m = p} NC_n(\{a\}) \tensor NC_{b_1}(X) \tensor \cdots \tensor NC_{b_q}(X)\tensor NC_m(\{b\}) \\
			&\to \bigoplus_{b_1 + \cdots + b_q = p}  NC_{b_1}(X) \tensor \cdots \tensor NC_{b_q}(X) =: N\mathscr{D}_{-p,q}.
			\end{split}
		\end{equation}
		We now check that the differentials on $N\Cobar$ and $N\mathscr{D}$ are compatible with $\phi$. Because all differentials of $NC_\bullet(\{a\})$ and $NC_\bullet(\{b\})$ are zero, the horizontal cobar differential formula \ref{eqn: twoside_cobar_differential} yields 
		\begin{equation*}
		\begin{split}
			\phi(d_H^{N\Cobar}(m \tensor c_1\tensor \cdots \tensor c_q\tensor n)) &= \phi(\sum_{i=1}^q (-1)^{\sigma(c_i)} m \tensor c_1 \tensor \cdots \tensor \partial_{NC}(c_i) \tensor \cdots \tensor c_q \tensor n) \\
			&= \sum_{i=1}^q (-1)^{\sigma(c_i)}  \psi_a(m\tensor c_1) \tensor \cdots \tensor \partial_{NC}(c_i) \tensor \cdots \tensor \psi_b(c_q\tensor n)  \\
			&= d_{N\mathscr{D}}(\phi(m\tensor c_1\tensor \cdots \tensor c_q \tensor n)).
		\end{split}
		\end{equation*}
		The vertical differential in both cases is given by the alternating sum of coface maps $d^i$. We will show that $\phi\circ d_{\Cobar}^i = d_{N\mathscr{D}}^i \circ \phi$ for all $i$. For $i=0$, the map $d_{N\mathscr{D}}^i: N\mathscr{D}^q_\bullet \to N\mathscr{D}^{q+1}_\bullet$ is induced by the following inclusion under the topological Alexander--Whitney map: 
		\[(x_1, ..., x_q)\to (a, x_1, ..., x_q).\]
		Naturality of Alexander--Whitney implies that the following diagram commutes:
			\[\begin{tikzcd}[sep=large]
				% https://q.uiver.app/#q=WzAsNixbMSwwLCJDX1xcYnVsbGV0KCpcXHRpbWVzIFhecSkiXSxbMiwwLCJDX1xcYnVsbGV0KFhee3ErMX0pIl0sWzEsMSwiQ19cXGJ1bGxldCgqKSBcXHRlbnNvciBDX1xcYnVsbGV0KFgpXntcXHRlbnNvciBxfSJdLFsyLDEsIkNfXFxidWxsZXQoWClee1xcdGVuc29yIHErMX0iXSxbMCwwLCJDX1xcYnVsbGV0KFhecSkiXSxbMCwxLCJDX1xcYnVsbGV0KFgpXntcXHRlbnNvciBxfSJdLFsyLDMsIkNfXFxidWxsZXQoXFxpb3RhX3gpXFx0ZW5zb3IgXFxpZF57XFx0ZW5zb3IgcX0iXSxbMCwyLCJcXEFXX3sqLCBYXnF9IiwyXSxbMCwxLCJDX1xcYnVsbGV0KFxcaW90YV94XFx0aW1lc1xcaWRfe1hecX0pIl0sWzEsMywiXFxBV197cSsxfSJdLFs0LDAsIlxcY29uZyJdLFs0LDUsIlxcQVdfcSIsMl0sWzUsMiwiXFxjb25nIl0sWzQsMSwiQ19cXGJ1bGxldChkXjApIiwwLHsiY3VydmUiOi01fV0sWzUsMywiXFxtYXRoc2Nye0R9KGReMCkiLDIseyJjdXJ2ZSI6NX1dXQ==
				{NC_\bullet(X^q)} & {NC_\bullet(\{a\}\times X^q)} & {NC_\bullet(X^{q+1})} \\
				{NC_\bullet(X)^{\tensor q}} & {NC_\bullet(\{a\}) \tensor NC_\bullet(X)^{\tensor q}} & {NC_\bullet(X)^{\tensor q+1}}
				\arrow["\cong", from=1-1, to=1-2]
				\arrow["{NC_\bullet(d^0)}", curve={height=-30pt}, from=1-1, to=1-3]
				\arrow["{\AW_q}"', from=1-1, to=2-1]
				\arrow["{NC_\bullet(\iota_a\times\id_{X^q})}", from=1-2, to=1-3]
				\arrow["{\AW_{\{a\}, X^q}}"', from=1-2, to=2-2]
				\arrow["{\AW_{q+1}}", from=1-3, to=2-3]
				\arrow["\cong", from=2-1, to=2-2]
				\arrow["{d^0_{\mathscr{D}}}"', curve={height=30pt}, from=2-1, to=2-3]
				\arrow["{NC_\bullet(\iota_a)\tensor \id^{\tensor q}}", from=2-2, to=2-3]
			\end{tikzcd}\]
		where $\iota_a$ is the inclusion $\{a\} \hookrightarrow X$. But now we make the identifications 
		\[NC_\bullet(\{a\})\tensor NC_\bullet(X)^{\tensor q}\cong NC_\bullet(X)^{\tensor q}, \text{ and } NC_\bullet(\iota_a) = \eta_a\]
		to conclude 
		\[ (\eta_a\tensor \id^{\tensor q}) \circ \AW_q = \AW_{q+1} \circ  \hspace{2pt}NC_\bullet(d^0)\]
		which implies $d_{N\mathscr{D}}^0 = \eta_a\tensor \id^{\tensor q}$. (By a slight abuse of notation we use $\eta_a$ to denote both the constant map at $\eta_a$ as well as the constant zero simplex at $a$). Then by naturality of the Alexander--Whitney map,
		\begin{equation*}
\begin{split}
    \phi(d^0_{\Cobar}(m \tensor c_1\tensor \cdots \tensor c_q\tensor n)) &= \phi(\Delta^R(m) \tensor c_1 \tensor \cdots \tensor c_q \tensor n) \\
    &= \phi\big(((\id\tensor \iota_a)\circ \AW(m)) \tensor c_1 \tensor \cdots \tensor c_q \tensor n \big) \\
    &= \phi(m \tensor \eta_a \tensor c_1 \tensor \cdots \tensor c_q \tensor n) \\
    &= \psi_a(m \tensor \eta_a) \tensor c_1 \tensor \cdots \tensor \psi_b(c_q \tensor n) \\
    &= \eta_a \tensor \psi_a(m \tensor c_1) \tensor \cdots \tensor \psi_b(c_q \tensor n) \\
    &= (\eta_a \tensor \id^{\tensor q})\big(\psi_a(m \tensor c_1) \tensor \cdots \tensor \psi_b(c_q \tensor n)\big) \\
    &= d^0_{N\mathscr{D}}(\phi(m \tensor c_1\tensor \cdots \tensor c_q \tensor n))
\end{split}
\end{equation*}
		The proof for $d^1, ..., d^q$ are analogous. Hence the restriction of $\phi$ to the truncated bicomplexes yields a quasi-isomorphism $N\Cobar^{\leq n} \to N\mathscr{D}^{\leq n} $, completing the proof.
	\end{proof}
\end{thm}
% \begin{rem}
% We have not yet used Adams' theorem. But what we have proven should suggest that the theorem is true for the following reason. We have exhibited a quasi-isomorphism between the two-sided cobar construction on $C_\bullet(X)$ and the totalization of $\mathscr{C}$, which is itself the bicomplex constructed by taking singular chains on the cosimplicial path space. So if we can somehow commute ``chains on a totalization'' and ``totalization of chains'' as in the right arrow in the following diagram
% \[\begin{tikzcd}[cramped]
% 	{\Cobar(C_\bullet(X))} && {\tot(\mathscr{C})=\tot(C_\bullet P^\bullet_{x,x}X)} \\
% 	{C_\bullet(\Omega_xX)} && {C_\bullet(\tot(P^\bullet_{x,x}X))}
% 	\arrow["\cong", from=1-1, to=1-3]
% 	\arrow["{\text{Adams' theorem}}"', dashed, from=1-1, to=2-1]
% 	\arrow["{?}", dashed, from=1-3, to=2-3]
% 	\arrow["\cong", from=2-1, to=2-3]
% \end{tikzcd}\]
% then we would have a version of Adams' theorem.
% \end{rem}
\begin{proof}[Proof of Theorem \ref{thm:looijenga}]
	By Theorem \ref{thm:two_sided_cobar_D}, Lemma \ref{lem:quasi_iso_CtoD}, and the discussion in the first subchapter, we have a quasi-isomorphism 
	\[\tot(\Cobar^{\leq n}) \xrightarrow{\ref{thm:two_sided_cobar_D}} \tot(\mathscr{D}^{\leq n}) \xrightarrow{\ref{lem:quasi_iso_CtoD}} \tot(\mathscr{C}^{\leq n}) \xrightarrow{ chapter } \tot(\mathscr{N}^{\leq n}).\]
	Therefore, on zeroth homology we have an isomorphism 
    \begin{equation}\label{eqn:cobar_iso_N}
		H_0\tot\Cobar^{\leq n}\to H_0\tot(\mathscr{N}^{\leq n})
	\end{equation}
	Now Theorem \ref{thm:stallings_generalization} yields an isomorphism 
	\begin{equation}\label{eqn:N_and_fundamental_gp}
		H_0\tot(\mathscr{N}^{\leq n}) \cong \Z\pi_X(a,b)/\Z\pi_X(a,b)\mathcal{I}^{n+1}.
	\end{equation}
    Combining Equations \ref{eqn:cobar_iso_N}, \ref{eqn:N_and_fundamental_gp} and applying Corollary \ref{cor:spectral_convergence}, we obtain a surjection
    \[\Z\pi_X(a,b)/ \Z\pi_X(a,b)\mathcal{I}^{n+1} \xrightarrow{\cong} H_0\tot\Cobar_{\leq n} \xrightarrow{\cong} H_0\tot_\bullet(\mathscr{N}^{\leq n})\twoheadrightarrow H_n(X^n, X(n)^x_x)\]
    whose kernel is isomorphic to $\Z$ if $a=b$ and trivial otherwise. It remains to see that if $a=b$, the kernel is generated by the constant path at $a$. 
\end{proof}
\subsection{Looking forward}
